\section*{Reviewer 3 --- Major Comments}
\reviewitem{R3.M1 --- Relation to submitted references [20] and [31]}{Identify inherited components and additional contributions; include the closest task-matched baseline.}{We reread both originals and revised the attribution. Iterative corridor construction, external penalties, and separate kinematic/separation penalties are acknowledged as antecedent ideas. The contribution list now focuses on the fixed common-time rest-to-rest formulation, affine QP structure, reachable-speed padding, endpoint moment correction with exact rollout, and the new hard-entry and acceptance-feedback procedure. The nonemptiness proof explains this local formulation; it is not presented as a general convergence theorem. Table 1 and the final Related Work paragraph give these distinctions. Hard box--iterative and Soft LIOM use the same UAV task, reference, corridors, reconstruction opportunities, and acceptance rules, providing controlled task-matched adaptations of the closest mechanisms. The revised comparison also includes task-matched ASCO, direct nonlinear programming, and sequential convex programming, with common endpoint, motion, initialization, and final-acceptance conditions. The new manuscript subsection ``Task-Matched Nonlinear and Sequential Convex Comparisons'' presents the paired case results. The ASCO risk-range update follows the original paper and archived implementation. SCP uses the original distance-norm tangent with sparse numeric OSQP assembly. The ten fixed tasks include ordinary crossings/passages, two winding layouts, and 15- and 20-UAV tasks. The table reports each task, its computation time, physical acceptance, and any budget overrun; no population success-rate ranking is inferred.}
\reviewitem{R3.M2 --- Physical difficulty}{Add denser teams, narrower passages, and multiple random instances.}{The new 60-task physical benchmark includes ten randomized crossing layouts at each of 10, 15, and 20 UAVs, plus five layouts at each of three passage/obstacle levels with ten UAVs. The workspace is $16\times16\times7$ m; passage openings are 3.2, 2.6, and 2.0 m. The design and seeds were saved before execution, and all outcomes are retained. These tests change physical congestion while preserving nominal corridor-growth settings. A further winding-counterflow study uses three alternating physical openings. In the initial three paired repetitions, the representative five-UAV task is accepted by the complete method and not by iterative hard RSC. The additional external-baseline study also records a hard-RSC acceptance of that task in 25.14 s. The original positive phase-I certificate applies to its initial corridors, and reconstruction can change their compatibility. A separate set of 12 previously untested realizations yields seven versus six accepted tasks. A physical-width/team/horizon map shows the additional recovery region, with every attempted outcome retained. \PhysicalEvidence}
\reviewitem{R3.M3 --- Equal reconstruction opportunities}{Give hard methods the same outer loop; the growth-cap sweep confounds iteration with softening.}{Hard box--iterative and Hard RSC--iterative now use the same reconstruction opportunities, reference-update rules, and static-guide policy. A usable numerical iterate can guide the next construction even when it is not accepted as a flight; otherwise the shared spatial-center/guide fallback applies. Every accepted motion passes identical checks. One-shot hard variants isolate the iteration restriction. Corridor-size changes are assessed as parameter sensitivity, separately from physical difficulty. In the new winding task, the initial hard RSC phase-I LP has matching positive primal and dual objective 0.21915. Removing only the relative rows or only the kinematic equalities makes that diagnostic LP feasible at zero violation. Thus the example identifies the coupled conflict while giving the hard baseline the same reconstruction opportunities. \HardIterationEvidence}
\reviewitem{R3.M4 --- Hard-first strategy and exact penalties}{Avoid multiple relaxed solves when an initial hard QP suffices; discuss finite exact-penalty thresholds.}{The proposed method begins with a bounded hard RSC-QP attempt. If its reconstructed motion passes checking, it returns immediately; otherwise the relaxed box-corridor loop follows. Soft LIOM removes this entry step. A further controlled study varies the entry allowance between 1, 3, and 5 s on ten tasks with identical references and a 30 s total budget. For C29, the 3 s allowance reduces total planning time from 19.27 to 3.63 s and QP calls from five to one. The 3 s setting is recommended for new runs; archived comparisons retain their explicitly labeled 1 s setting. Table~\ref{tab:entry-budget-response} reproduces the full ablation.  Section 6.2 discusses the linear slack penalty: under the usual convex-QP duality assumptions, a sufficiently large weight relative to optimal multipliers can be exact for those inequalities. A quadratic defect penalty generally remains inexact at finite weight. Caps therefore address conditioning, while measured recovery and final checking determine acceptance. No universal finite exactness threshold is asserted for both terms. \HardEntryEvidence}
\begin{table}[htbp]

\caption{Hard-entry budget ablation. Each cell gives total planning time in seconds and QP calls. All three settings share the recorded reference and a 30 s total budget, with rotated execution order. $\times$ marks no accepted motion and $\dagger$ marks late acceptance.\label{tab:entry-budget-response}}
\small\begin{tabularx}{\textwidth}{Xrrr}\toprule
Task & 1 s entry: time / QPs & \textbf{3 s entry: time / QPs} & 5 s entry: time / QPs\\\midrule
C01 & 1.75 / 1 & 0.95 / 1 & 0.88 / 1\\
C03 & 1.35 / 1 & 1.44 / 1 & 1.46 / 1\\
C05 & 2.84 / 1 & 2.91 / 1 & 2.89 / 1\\
C25 & 1.20 / 1 & 1.18 / 1 & 1.17 / 1\\
C27 & 24.08 / 6 & 26.15 / 6 & 28.29 / 6\\
C29 & 19.27 / 5 & 3.63 / 1 & 3.65 / 1\\
Winding A & 28.70 / 5 & 30.00 / 5 & 31.33$^{\dagger}$ / 5\\
Winding B & 30.48$^{\dagger}$ / 6 & $\times$\,30.07 / 5 & $\times$\,30.04 / 5\\
15-UAV & $\times$\,30.08 / 7 & 29.09 / 6 & $\times$\,30.66 / 6\\
20-UAV & 28.37 / 6 & 27.04 / 6 & 30.35$^{\dagger}$ / 6\\
\bottomrule\end{tabularx}
\end{table}

\reviewitem{R3.M5 --- Reconstructed-motion acceptance}{Clarify the candidate/output gap, add feedback after failed checks, and normalize recovery.}{The final output comes from corrected controls and exact integration. Its independent checker examines every obstacle and UAV pair, endpoints, workspace, and Euclidean speed/acceleration limits. A numerically recovered candidate that fails this check triggers a tighter normalized defect tolerance and a capped defect-weight increase, followed by an exact-kinematics, zero-slack attempt in the current boxes. If that also fails, reconstruction continues within budget. The normalized defect is $\bar D=(h_0/h)^2D/(NK)$; slack is averaged over active pair--node restrictions. These are numerical criteria, not substitutes for physical acceptance. The acceptance-feedback ablation evaluates this combined procedure and records rejected-check reasons. \AcceptanceEvidence}
\reviewitem{R3.M6 --- Motion-envelope conservatism}{Replace or justify the inscribed cube and use physical Euclidean bounds for final checks.}{We adopted a richer polyhedral inner approximation to retain linear QP constraints. The nominal convex hull of six unit-axis directions and eight normalized cube vertices has 24 facets. It contains the old cube, reaches full axial authority, and lies inside the Euclidean ball. Scaling gives velocity and acceleration limits. Final checks now use Euclidean norms. The original cube and a 26-vertex, 48-facet envelope are compared with the nominal envelope. \EnvelopeEvidence}
\reviewitem{R3.M7 --- Theory and termination statistics}{Nonemptiness does not prove convergence; report rounds and failure branches.}{The proposition is explicitly titled ``Nonemptiness of the relaxed fixed-corridor subproblem.'' It assumes nonempty, endpoint-compatible spatial boxes and establishes a local feasible point. Numerical recovery and final physical acceptance are defined separately. No general outer-convergence or bounded-round success claim is made. Tables give outer rounds and QP calls. Records retain initialization/construction failures, numerical nonrecovery, rejected reconstructed motions, stalled hard reconstruction, and budget exits. \FailureEvidence}
\reviewitem{R3.M8 --- QP solver and timing decomposition}{Use a dedicated solver, compare matched QPs, and report stage and scale effects.}{Sparse numeric QPs now use OSQP. End-to-end timing includes reference generation, corridors, assembly, solver setup/data exchange, iterations, endpoint reconstruction, and verification. CPU time includes MATLAB and the worker. A separate identical-matrix OSQP audit compares candidate-readiness stopping with ordinary solver termination at two tolerances. All planning QPs use OSQP. The additional ASCO and direct-NLP comparisons use Ipopt with ma57 for their nonlinear collision constraints. Both baselines are rerun on all ten fixed tasks with the same inputs and budgets. The $N$--$K$ study and stage plot address scale effects. \SolverEvidence}
\reviewitem{R3.M9 --- Physical realism}{Add higher-fidelity tracking or state the model scope.}{MuJoCo simulations include six-degree-of-freedom rigid-body dynamics, four bounded rotor forces, motor lag, gravity, drag, and collision geometry. Position/velocity feedback drives a rotation-matrix attitude controller. Fixed controller settings are used under nominal and perturbed mass, lag, and wind. Terminal error is measured at $T_f$; the logged settling interval is separate. The Conclusion identifies the static-map setting and additional modeling needed for downwash and perception uncertainty. The complete controller and aircraft simulation execute in software; no flight-control hardware is connected. \TrackingEvidence}
\reviewitem{R3.M10 --- Initializer effects}{Assess SFM sensitivity/failure and separate reference geometry from later guide availability.}{The initializer study crosses SFM, static-guide, and straight geometry with guide availability during relocation. The guide-only fallback after SFM/timing failure is explicitly recorded. The sensitivity study varies repulsion, decay length, look-ahead, response time, and tangential force independently. All initialization failures and fallback reasons remain in the saved records. \InitializationEvidence}
